\documentclass[10pt,a4paper]{amsart}
\usepackage{amsmath,amssymb,amsthm,graphicx}
\usepackage{verbatim}
\newenvironment{reference}{\comment}{\endcomment}
%\newenvironment{reference}{}{}
\newenvironment{slogan}{\comment}{\endcomment}
\newenvironment{history}{\comment}{\endcomment}

% For commutative diagrams we use Xy-pic
\usepackage[all]{xy}

% We use 2cell for 2-commutative diagrams.
\xyoption{2cell}
% Theorem environments.
%
\theoremstyle{plain}
\newtheorem{theorem}[subsection]{Theorem}
\newtheorem{proposition}[subsection]{Proposition}
\newtheorem{lemma}[subsection]{Lemma}

\theoremstyle{definition}
\newtheorem{definition}[subsection]{Definition}
\newtheorem{example}[subsection]{Example}
\newtheorem{exercise}[subsection]{Exercise}
\newtheorem{situation}[subsection]{Situation}

\theoremstyle{remark}
\newtheorem{remark}[subsection]{Remark}
\newtheorem{remarks}[subsection]{Remarks}

\numberwithin{equation}{subsection}

% Macros
%
\def\lim{\mathop{\mathrm{lim}}\nolimits}
\def\colim{\mathop{\mathrm{colim}}\nolimits}
\def\Spec{\mathop{\mathrm{Spec}}}
\def\Hom{\mathop{\mathrm{Hom}}\nolimits}
\def\Ext{\mathop{\mathrm{Ext}}\nolimits}
\def\SheafHom{\mathop{\mathcal{H}\!\mathit{om}}\nolimits}
\def\SheafExt{\mathop{\mathcal{E}\!\mathit{xt}}\nolimits}
\def\Sch{\mathit{Sch}}
\def\Mor{\mathop{\mathrm{Mor}}\nolimits}
\def\Ob{\mathop{\mathrm{Ob}}\nolimits}
\def\Sh{\mathop{\mathit{Sh}}\nolimits}
\def\NL{\mathop{N\!L}\nolimits}
\def\CH{\mathop{\mathrm{CH}}\nolimits}
\def\proetale{{pro\text{-}\acute{e}tale}}
\def\etale{{\acute{e}tale}}
\def\QCoh{\mathit{QCoh}}
\def\Ker{\mathop{\mathrm{Ker}}}
\def\Im{\mathop{\mathrm{Im}}}
\def\Coker{\mathop{\mathrm{Coker}}}
\def\Coim{\mathop{\mathrm{Coim}}}

% Boxtimes
%
\DeclareMathSymbol{\boxtimes}{\mathbin}{AMSa}{"02}

%
% Macros for moduli stacks/spaces
%
\def\QCohstack{\mathcal{QC}\!\mathit{oh}}
\def\Cohstack{\mathcal{C}\!\mathit{oh}}
\def\Spacesstack{\mathcal{S}\!\mathit{paces}}
\def\Quotfunctor{\mathrm{Quot}}
\def\Hilbfunctor{\mathrm{Hilb}}
\def\Curvesstack{\mathcal{C}\!\mathit{urves}}
\def\Polarizedstack{\mathcal{P}\!\mathit{olarized}}
\def\Complexesstack{\mathcal{C}\!\mathit{omplexes}}
% \Pic is the operator that assigns to X its picard group, usage \Pic(X)
% \Picardstack_{X/B} denotes the Picard stack of X over B
% \Picardfunctor_{X/B} denotes the Picard functor of X over B
\def\Pic{\mathop{\mathrm{Pic}}\nolimits}
\def\Picardstack{\mathcal{P}\!\mathit{ic}}
\def\Picardfunctor{\mathrm{Pic}}
\def\Deformationcategory{\mathcal{D}\!\mathit{ef}}
\usepackage{cleveref}
\crefname{originallemma}{original lemma}{original lemmas}
\Crefname{originallemma}{Original lemma}{Original lemmas}
\makeatletter
\newlabel{lemma-ML-limit-nonempty}{{0597}{source}{}{sourcepointer-1}{}}
\newlabel{lemma-ML-limit-nonempty@cref}{{[originallemma][0][]0597}{[1][0][]source}{}{}{}}
\makeatother
\title{A Mittag-Leffler kernel makes the limit of a short exact sequence of countable inverse systems exact}
\author{The Stacks Project Authors}
\date{stacks-project:algebra.tex:local-snapshot-20261012; GNU FDL 1.2 or later}
\begin{document}
\maketitle
\noindent\textit{Editorial note.} All delivered mathematical text is the author's own, copied from the pinned Stacks Project chapter file algebra.tex; no statement or proof step was rewritten, reordered or removed.\par\medskip
\noindent\textit{Editorial note.} Preamble provenance, as required for this source: the frozen snapshot stages the chapter's own preamble file beside it (preamble.tex, 217 lines, 8116 bytes, the same shared Stacks-wide preamble the published chow and etale-cohomology cards use); the first line of the pinned chapter file is \textbackslash{}input\{preamble\}, so no preamble text lives inside the chapter itself. The card takes that file's lines 8-18 (the verbatim package and the three \textbackslash{}comment-based environments reference/slogan/history at lines 9-12, section 13.2, together with the Xy-pic setup) and its lines 154 to the runtime end of the file (the native theorem declarations and mathematical macros). Not one character was changed, no package was added, no macro was defined, and the excerpt serves only as preamble, never as a span coordinate.\par\medskip
\noindent\textit{Editorial note.} The same-chapter pointers are located at their \textbackslash{}label lines inside the pinned chapter file; the cross-chapter pointers are located in the project's own tag list staged in the snapshot (sources/stacks-chow/tags.txt), at the verbatim lines carrying their canonical tags.\par\medskip
\noindent\textit{Editorial note.} The counted proof refers to a same-chapter statement by label; it is recorded as a navigation stub pointed at its label line.\par\medskip
\noindent\textit{Editorial note.} The counted spans contain no cite, so the card carries no bibliography.\par\medskip
\noindent\textit{Allowed standard prerequisites (editorial).}\begin{itemize}
\item Left exactness of inverse limits of directed systems. Used at the start of the counted proof (line 21001). Reduces the claim to surjectivity at the last term.
\item The Mittag-Leffler property as bounded image stabilisation. Used throughout the counted proof. Lets the preimages of a compatible family be chosen compatibly along the system.
\end{itemize}\medskip
\ifdefined\hypertarget\hypertarget{sourcepointer-1}{}\fi
\noindent\textit{Source reference.} Original lemma 0597: sources/stacks-algebra/pinned-original-source.tex, source lines 20965--20965. Navigation pointer for \textbackslash{}ref\{lemma-ML-limit-nonempty\} used inside the counted span (line 21007); the argument that consumes it is retained in the span.\par\medskip
\begin{lemma}
\label{lemma-ML-exact-sequence}
Let
$$
0 \to A_i \xrightarrow{f_i} B_i \xrightarrow{g_i} C_i \to 0
$$
be an exact sequence of directed inverse systems of abelian groups over $I$.
Suppose $I$ is countable.  If $(A_i)$ is Mittag-Leffler, then
$$
0 \to \lim A_i \to \lim B_i \to \lim C_i\to 0
$$
is exact.
\end{lemma}

\begin{proof}
Taking limits of directed inverse systems is left exact, hence we only need to
prove surjectivity of $\lim B_i \to \lim C_i$.  So let $(c_i) \in \lim
C_i$.  For each $i \in I$, let $E_i = g_i^{-1}(c_i)$, which is nonempty since
$g_i: B_i \to C_i$ is surjective. The system of maps $\varphi_{ji}: B_j
\to B_i$ for $(B_i)$ restrict to maps $E_j \to E_i$ which
make $(E_i)$ into an inverse system of nonempty sets.  It is enough to show
that $(E_i)$ is Mittag-Leffler. For then Lemma \ref{lemma-ML-limit-nonempty}
would show $\lim E_i$ is nonempty, and taking any element of $\lim E_i$ would
give an element of $\lim B_i$ mapping to $(c_i)$.

\medskip\noindent
By the injection $f_i: A_i \to B_i$ we will regard $A_i$ as a subset of
$B_i$.  Since $(A_i)$ is Mittag-Leffler, if $i \in I$ then there exists $j \geq
i$ such that $\varphi_{ki}(A_k) = \varphi_{ji}(A_j)$ for $k \geq j$.  We claim
that also $\varphi_{ki}(E_k) = \varphi_{ji}(E_j)$ for $k \geq j$.  Always
$\varphi_{ki}(E_k) \subset \varphi_{ji}(E_j)$ for $k \geq j$.  For the reverse
inclusion let $e_j \in E_j$, and we need to find $x_k \in E_k$ such that
$\varphi_{ki}(x_k) = \varphi_{ji}(e_j)$.  Let $e'_k \in E_k$ be any element,
and set $e'_j = \varphi_{kj}(e'_k)$.  Then $g_j(e_j - e'_j) = c_j - c_j = 0$,
hence $e_j - e'_j = a_j \in A_j$.  Since $\varphi_{ki}(A_k) =
\varphi_{ji}(A_j)$, there exists $a_k \in A_k$ such that $\varphi_{ki}(a_k) =
\varphi_{ji}(a_j)$.  Hence
$$
\varphi_{ki}(e'_k + a_k) = \varphi_{ji}(e'_j) + \varphi_{ji}(a_j) =
\varphi_{ji}(e_j),
$$
so we can take $x_k = e'_k + a_k$.
\end{proof}

\end{document}
